% Transcription — fonds Alexandre Grothendieck, shelfmark 148, batch 1 (pages 1-20).
% Established from the facsimile archives/batches/148/batch-01.pdf, cut from a
% copy of 148.pdf fetched through the project's relay
% (https://grothendieck-relay.onrender.com/source/148.pdf); PDF page k is the
% archivists' page k, no cover sheet. Per the skill transcribe-grothendieck.
% The folder is seventy-one pages in four batches; this is the first.
% Batches 3 and 4 were read for notation; batch 2 is written in parallel.
%
% Pass: Opus 5.5 (claude-opus-5-5), 2026-09-26 — first pass, unchecked against the pages by a human.
%
% Pages skipped: 5, 7, 9, 11, 13, 15 — the versos of his leaves 4-14,
% typed USTL « avis de soutenance » circulars (dated 14, 18 and 20 October
% 1983), with nothing of his; skipped by the settled rule. Page 16 is a
% cover: blank listing paper bearing only « M_{0 4} » and « M_{1,1} » in his
% hand at top right; it takes no \page and his words become the section
% heading, with a \note. Pages summarised: none (no typescript). Page 20 (back
% of a listing) carries only a pencil sketch of his and is kept with a
% \note describing it.
%
% Dates: page 1 is dated in his hand « (29.11.83) » in the left margin; it is
% recorded in a \note there. The verso circulars (Oct. 1983) are not his
% and are not recorded in the file. \dating is the catalogue's « 1983 »,
% verbatim, with no group sentence.
%
% His own pagination: pages 1-3 carry no number of his. Pages 4, 6, 8, 10,
% 12, 14 form a run he numbers 1 to 6 at top centre (« Courbes exceptionnelles
% de type (1,1) », A) square type, B) hexagonal type), recorded in \note{}s at
% each page. Pages 17-20 (listing paper) are unnumbered. This run is distinct
% from the 1-21 run that starts on page 43 (batches 3-4).
%
% What the batch is about. Pages 1-3 (dated 29.11.83): for a compact oriented
% surface X with toric structures on its boundary circles, the group
% SRT(X) (automorphisms modulo the identity component of those that are the
% identity on the boundary), its three-step filtration by S_*T^!(X), U^I,
% S_I (I = pi_0(boundary)), a canonical subgroup R^I ⊃ Z^I with
% SRT^! = R^I · S_*T^!, the extension 1 -> U^I -> SUT(X) -> T(X) -> 1,
% the splittings attached to a marking s ∈ ∏ T_i and the embedding
% S_*T(X, s) -> SRT(X), and the multimarked generalisation. Pages 4-14 (his
% 1-6): the two exceptional elliptic curves of type (1,1) — the square
% (C/Z[i], torsors under Z/4 = mu_4) and the hexagonal (C/Z[zeta], torsors
% under Z/6 = mu_6) — their fixed points of order 2 and 3, their canonical
% cell decompositions and « doubled » decompositions, the quotients by ±1
% (sphere subdivided into two squares; tetrahedral structure with a fixed
% vertex). Pages 16-20 (listing paper, coloured inks and pencil): M_{0,4}
% and M_{1,1} — the truncated tetrahedron, its barycentric subdivision,
% counts s_0, s_1, s_2, the four types of base points (square, tetrahedral,
% « pantalonnic » structures) and of elementary paths a)-g); on page 19 the
% configurations Q_0, Q_1, Q_infty, relations rho^3 = sigma^2 = 1,
% rho sigma = l, and a diagram of exact sequences linking T_{0,4}, T'_{0,4},
% T_{1,1} = Sl(2,Z) to S_4 = Aff(2,Z/2), S_3 = Sl(2,Z/2), Sl(2,Z/4)/±.
%
% Notation, aligned with batches 3-4: S_* with the star under the S; his
% gothic S as \mathfrak{S}; the circle group, a double-struck U, as
% \mathbb{U}; his double-struck El as \mathbb{E}\mathrm{l}; contracted
% products as \wedge with the group as superscript; T (roman) for the
% groups of pages 1-3, script \mathcal{T} for those of page 19 as he draws
% them.
%
% Hardest pages: 1 (oblique margin note, mostly illegible; the ! exponents);
% 17 (crowded listing page with drawings, oblique margin note, footnote at
% the foot); 19 (pencil, faint). Readings to check first: the exponent
% « ! » on page 1, « i-fix » on page 6, « régiments » on page 10, the
% genus words and the margin on page 17, the first term of the last row
% of the page 19 diagram.
%
\documentclass[11pt,a4paper]{article}
% Le préambule commun à toutes les transcriptions du fonds Grothendieck.
%
% Il tient en une idée : **l'incertitude se compose, elle ne se résout pas.**
% Une transcription qui lisse un mot douteux a détruit la seule chose pour
% laquelle on la faisait — un lecteur doit pouvoir savoir, à chaque endroit,
% ce qui était sur le feuillet et ce qui a été ajouté. Les six macros qui
% suivent sont donc l'appareil critique, et non de la décoration.
%
% Les mêmes commandes sont reconnues par `scripts/render.mjs`, qui produit la
% vue de lecture du site. Une macro ajoutée ici doit l'être aussi là-bas,
% faute de quoi le rendu échouera bruyamment — ce qui est voulu : un rendu
% silencieusement faux est pire qu'un rendu absent.

\ProvidesPackage{grothendieck}[2026/08/08 Transcriptions du fonds Grothendieck]

\usepackage{amsmath,amssymb,amsthm}
\usepackage{mathrsfs}
\usepackage{stmaryrd}
% Les diagrammes commutatifs sont partout dans ce fonds : c'est la langue
% même de la géométrie algébrique de Grothendieck, et une transcription qui
% les rendrait en prose ne transcrirait rien.
\usepackage{tikz-cd}
% Français par défaut — c'est la langue des feuillets — mais l'anglais est
% chargé aussi : la traduction et le résumé sont des documents anglais, et la
% typographie française (espaces avant les deux-points) y serait une faute.
% Ces documents basculent par \selectlanguage{english} en tête de corps.
\usepackage[english,french]{babel}
\usepackage{fontspec}
\usepackage{geometry}
\geometry{a4paper,margin=25mm,marginparwidth=28mm}
\usepackage{marginnote}
\usepackage{xcolor}
% ulem, not soul. `soul`'s \st and \ul are built on a character-by-character
% scan that XeLaTeX's font machinery breaks: the document fails with an
% undefined control sequence rather than degrading. ulem does the same two jobs
% and is Unicode-engine safe. `normalem` stops it hijacking \emph, which the
% transcriptions use for Grothendieck's own emphasis.
\usepackage[normalem]{ulem}
\usepackage{hyperref}
\hypersetup{colorlinks=true,linkcolor=black,urlcolor=black,pdfborder={0 0 0}}

% --- Métadonnées du lot -----------------------------------------------------
% Reprises telles quelles par le script de rendu, qui les affiche en tête de
% la vue de lecture. Elles fixent ce que le fichier prétend couvrir : sans
% elles, une transcription flotte sans point d'attache dans un fonds de seize
% mille feuillets.

\newcommand{\folder}[1]{\gdef\grothFolder{#1}}
\newcommand{\batch}[1]{\gdef\grothBatch{#1}}
\newcommand{\leaves}[2]{\gdef\grothFirst{#1}\gdef\grothLast{#2}}
\newcommand{\foldertitle}[1]{\gdef\grothFoldertitle{#1}}
\gdef\grothFolder{?}\gdef\grothBatch{?}\gdef\grothFirst{?}\gdef\grothLast{?}\gdef\grothFoldertitle{}

% --- Le résumé --------------------------------------------------------------
%
% Il ouvre la lecture modernisée, sous le titre « Résumé », et il est composé
% en petit. Ce n'est pas un ornement : c'est le seul passage du document qui
% ne soit pas des mathématiques, et un lecteur doit pouvoir le reconnaître
% avant de l'avoir lu — puis le sauter s'il connaît déjà le sujet, ou s'y
% arrêter s'il ne le connaît pas. Le filet qui le suit marque la reprise du
% corps.

\newenvironment{resume}
  {\small\setlength{\parskip}{0.45em}}
  {\par\medskip\hrule\medskip}

% --- Mots-clés ---------------------------------------------------------------
%
% En anglais, en clôture du résumé de la lecture modernisée : le vocabulaire
% moderne sous lequel chercher ce que les feuillets construisent. Le manifeste
% du site (`npm run manifest`) les extrait pour étiqueter la cote entière —
% cette ligne est la source unique des tags, il n'y a pas de fichier à côté.
\newcommand{\keywords}[1]{\par\smallskip\noindent
  {\footnotesize\textsc{Keywords} --- \textit{#1}}\par}

% --- L'appareil critique ----------------------------------------------------

% Le numéro de feuillet, en marge. C'est l'ancre qui permet de régler un
% désaccord sur une formule en regardant *un* feuillet plutôt qu'une chemise
% de six cents. Le site s'en sert aussi pour tourner les pages du fac-similé
% quand on fait défiler la transcription.
\newcommand{\leaf}[1]{%
  \marginnote{\footnotesize\color{gray!70}\textsf{#1}}%
  \label{leaf:#1}\ignorespaces}

% Illisible : on ne devine pas. Un mot inventé qui se lit comme les autres est
% le pire résultat possible.
\newcommand{\ill}{\textcolor{red!60!black}{[\,\ldots\,]}}

% Lecture douteuse : le mot est donné, et signalé comme incertain.
\newcommand{\uncertain}[1]{\uline{#1}}

% Ajout de l'éditeur — une lettre manquante, un mot sous-entendu.
\newcommand{\add}[1]{\textcolor{blue!50!black}{[#1]}}

% Biffé par Grothendieck lui-même. Ce qu'il a rayé fait partie du document :
% une rature dit souvent où la pensée a changé de direction.
\newcommand{\struck}[1]{\sout{#1}}

% Note du transcripteur, sur l'état matériel du feuillet.
\newcommand{\note}[1]{\par\smallskip{\small\color{gray!60!black}\textit{[#1]}}\par\smallskip}

% Note marginale de Grothendieck, distincte du corps du texte.
\newcommand{\marginal}[1]{\par\smallskip\noindent\hspace{1em}%
  {\small\color{gray!60!black}#1}\par\smallskip}

% --- Titre ------------------------------------------------------------------

\AtBeginDocument{%
  \begin{center}
    {\large\bfseries Fonds Alexandre Grothendieck --- cote n\textsuperscript{o}~\grothFolder}\\[2pt]
    {\itshape\grothFoldertitle}\\[4pt]
    {\small feuillets \grothFirst--\grothLast\ (lot \grothBatch)}\\[2pt]
    {\footnotesize\color{gray!60!black}Transcription établie d'après le fac-similé numérisé
     par l'Université de Montpellier.\\
     Les lectures incertaines sont soulignées, les illisibles notées [\,\ldots\,],
     les ajouts entre crochets.}
  \end{center}
  \vspace{1em}}

\endinput


\folder{148}
\batch{1}
\pages{1}{20}
\dating{1983}
\watermark{Édition de démonstration}
\foldertitle{Teichmüller (1983) : notes manuscrites (1983)}

\begin{document}
\section{Groupes de Teichmüller à structures toriques}

\page{1}
\marginal{(29.11.83)}\note{date de sa main, écrite en oblique dans la
marge gauche, en face de la première ligne}

Soit $X$ surf.\ à bord compacte orientée, à structure torique donnée sur
les comp.\ connexes du bord\marginal{(i.e.\ torseur sous
$\mathbb{U} = \lbrace z \in \mathbf{C} \mid |z| = 1 \rbrace$)}\note{en
interligne au-dessus de la fin de la phrase, rattaché par un trait à
« structure torique »}. On définit
\[
SRT(X) =
\]
quotient du groupe des autom.\ de $X$ (\uncertain{respectant}
l'orientation et les structures toriques) par la comp.\ neutre du
sous-groupe $\mathrm{Aut}^{!}(\vec{X})$ \struck{\ill{}} de
$\mathrm{Aut}$ qui induit l'identité sur $\partial X$.

On a ainsi une \uncertain{filtration} de $SRT(X)$ en trois étages
\[
S_{*}T^{!}(X) \qquad \mathbb{U}^{I} \qquad \mathfrak{S}_{I}
\qquad\qquad I = \pi_0(\partial X)
\]
\note{sous $S_{*}T^{!}(X)$, en petit : « $\mathrm{Aut}^{!}(\vec{X})$/comp.\
neutre ». La lettre gothique qu'il écrit pour le groupe symétrique est
rendue $\mathfrak{S}$ dans tout le lot ; l'étoile de $S_{*}$ est écrite
sous le $S$, comme dans les lots 3 et 4}
ou encore
\[
\mathbf{Z}^{I} \qquad T(X) \qquad \mathbb{U}^{I} \qquad \mathfrak{S}_{I}
\]

En fait, on définit canon.\ un sous-groupe $\mathbf{R}^{I}$\marginal{droit}
\note{« droit » au-dessus de $\mathbf{R}^{I}$, rattaché par une accolade}
de $SRT(X)$ contenant $\mathbf{Z}^{I}$, de telle façon qu'on ait
\[
SRT^{!}(X) = \mathbf{R}^{I} \cdot S_{*}T^{!}(X)
\]
\note{sous $SRT^{!}(X)$, une accolade et « noyau de
$SRT^{!}(X) \to \mathfrak{S}^{!}_{I}$ » ; les signes $!$ de cette ligne
sont \uncertain{lus} tels quels}

\note{à gauche, un petit diagramme : $\mathbf{R}^{I}$ en haut, une flèche
d'inclusion, un trait biffé, une flèche verticale marquée
« \uncertain{cocart} », et un point ; puis la ligne ci-dessous}
\[
\mathbf{Z}^{I} \hookrightarrow S_{*}T^{!} \subset SRT^{!} \subset SRT
\qquad \text{i.e.\ on a}
\]

\marginal{\emph{NB} $SUT(X) \simeq T(X) \times_{\mathfrak{S}_I}
\mathrm{Aut}((T_i)_{i\in I})$}\note{en oblique, à droite}

\begin{tikzcd}[column sep=small, row sep=small]
 & SRT(X) \arrow[d, no head, "\mathfrak{S}_I"] & \\
 & \mathbf{R}^{I}\cdot S_{*}T^{!} \arrow[dl, no head, "T^{!}"'] \arrow[dr, no head, "\mathbb{U}^{I}"] & \\
\mathbf{R}^{I} \arrow[dr, no head, "\mathbb{U}^{I}"'] & & S_{*}T^{!}(X) \arrow[dl, no head, "T^{!}(X)"] \\
 & \mathbf{Z}^{I} &
\end{tikzcd}

\note{treillis de sous-groupes dessiné au bas de la page, traits sans
pointe ; les quotients sont écrits à côté des traits et placés ici en
étiquettes. À droite de $SRT(X)$, une accolade :
« $\mathrm{Aut}((T_i)_{i\in I})$ »}

d'où
\[
SRT/\mathbf{Z}^{I} \overset{\mathrm{def}}{=} SUT(X)
\quad\big\rbrace\ \mathfrak{S}_I
\]
\note{sous $SUT(X)$, une flèche vers le bas et « $T^{!}(X) \times
\mathbb{U}^{I}$ », surmonté d'une accolade}
\[
(*)\qquad 1 \to \mathbb{U}^{I} \to SUT(X) \to T(X) \to 1
\]

\marginal{Jusqu'ici on a fait intervenir $S_{*}T^{!}(X)$, i.e.\
\uncertain{l'exclusion} \ill{} de la définition $T(X)$ \ill{}
$S_{*}T(X)$, str.\ \ill{}. \ill{}}\note{note écrite en oblique dans la
marge gauche, lue sur l'image tournée, en grande partie illisible}

De plus, chaque fois qu'on se donne un \uncertain{marquage}, i.e.\ un
pt de
\[
\underline{s} = (s_i) \in \prod_{i \in I = \pi_0(\partial X)} T_i
\]
\note{après le produit, un mot biffé illisible}
($T_i$ = comp.\ connexe (torique) de $\partial X$ d'indice $i$), on trouve
un plongement
\[
S_{*}T^{!}(X, \underline{s}) \hookrightarrow SRT(X)
\]
tel qu'on \ill{} \uncertain{en déduire}
\[
S_{*}T^{!}(X, \underline{s}) \wedge^{\mathbf{Z}^{I}} \mathbf{R}^{I}
\simeq SRT(X)
\]
\note{l'exposant de $T$ dans ces deux formules est un signe surchargé,
lu $!$ sans certitude ; page 2 il écrit $S_{*}T(X,\underline{s})$ sans
exposant. Le produit contracté est écrit $\wedge$ avec $\mathbf{Z}^{I}$
en indice sous le signe}

\page{2}
Une façon de le voir est de dire que la donnée de $\underline{s}$
définit un scindage $\varphi_{\underline{s}}$ de l'extension $(*)$
ci-dessus, et que $S_{*}T(X, \underline{s})$ s'identifie au ss-groupe
de $SRT(X)$ image inverse de
\[
\varphi_{\underline{s}}(T(X)) \subset SUT(X)
\]
par l'homom.\ \uncertain{canon.}
\[
SRT(X) \to SUT(X)
\]
de passage au quotient\note{le premier groupe est écrit comme
« $SUR(X)$ » ; lu $SRT(X)$ d'après la page 1}. Le changement de
$\underline{s}$ en $\underline{s}' = u \cdot \underline{s}$ transforme
$\varphi_{\underline{s}}$ en
\[
\varphi_{\underline{s}'} = \mathrm{int}(u) \circ \varphi_{\underline{s}}
\]
donc (si $\tilde{u} \in \mathbf{R}^{I}$ relève $u$), on trouve
\[
S_{*}T(X, \underline{s}') = \mathrm{int}(\tilde{u})\, S_{*}T(X, \underline{s}).
\]

En termes de $\overline{X}$, la définition de $S_{*}T(X,\underline{s})$
dépend du choix $\underline{s}$ \emph{non canonique}, par contre celle de
$S_{*}RT(X)$ ne dépend que d'un choix canonique, celui de la famille des
structures toriques $(T_i)_{i\in I}$. De plus, $\prod T_i$ comme torseur
sous $\mathbb{U}^{I}$ se réalise \marginal{(par \uncertain{fidélité}
\ill{})} comme un \ill{} de scindages de l'ext.\ $(*)$ (via
\ill{} par $\mathbb{U}^{I}$) [\uncertain{déduite} de $SRT(X)$ sur les
structures \ill{}] \uncertain{savoir} les \uncertain{relèvements}
coïncidant sur $\mathbf{Z}^{I}$.

\begin{tikzcd}[column sep=small, row sep=small]
\mathbf{R}^{I} \arrow[dr] & \\
 & SRT(X) \\
S_{*}T^{!}(X) \arrow[ur] &
\end{tikzcd}

\page{3}
Ceci dit, si on se donne un \emph{multimarquage} de $\partial X$, savoir
pour chaque $i \in I$ un $n_i \in \mathbf{N}^{*}$ et un sous-torseur
\marginal{\uncertain{élém.} $n_i$}\note{en interligne au-dessus de
« torseur »}
\[
T_i(n_i) \subset T_i,
\]
\struck{de groupe} i.e.\ un $s_i \in T_i^{\wedge n_i}$, d'où
\struck{$\underline{s}$}
\[
\underline{s} \in \prod_{i\in I} T_i^{\wedge n_i},
\]
alors le groupe de Teichmüller multimarqué $S_{*}T(X, \underline{s})$,
généralisant le précédent, est plongé dans $SRT(X)$ de façon naturelle,
en prenant un $s_i$ dans chaque $T_i(n_i)$, d'où
\[
\underline{s}^{0} \in \prod T_i \quad\text{et}\quad
S_{*}T(X, \underline{s}^{0}) \subset SRT(X),
\]
et on prend le ss-groupe
\[
S_{*}T(X, \underline{s}) = S_{*}T(X, \underline{s}^{0}) \cdot
\prod_{i\in I} \tfrac{1}{n_i}\mathbf{Z}_i
\qquad\qquad \mathbf{Z}^{I} \subset SRT(X).
\]
où $\mathbf{Z}_i$ est le $i$-ème facteur de $\mathbf{Z}^{I}$.
\note{l'exposant de $T_i$, lu $\wedge n_i$, est écrit petit et haut ; le
reste de la page est blanc}

\section{Courbes exceptionnelles de type (1,1)}

\page{4}
\note{p.~1 de l'auteur. Ici commence une suite paginée par lui, en haut
au centre, sur les rectos seuls : sa p.~1 est la page 4, sa p.~2 la page
6, sa p.~3 la page 8 (voir plus bas pour la suite). Le titre de section
est le sien}

\emph{Courbes} \add{exceptionnelles}\note{« exceptionnelles » en
interligne au-dessus de « de type »} \emph{de type} (1,1) (sur
$\mathbf{C}$)

A) \emph{Courbe} \add{ell.}\ \emph{anharmonique de type carré}

Cat.\ de ces courbes équivalente : celle des carrés combinatoires
orientés, i.e.\ des torseurs sous
\[
\mathbf{Z}/4\mathbf{Z} \simeq \mu_4(\mathbf{C}) \simeq
\]
groupe des unités algébriques de $\mathbf{Q}[i]$, i.e.\ groupe des unités
de $\mathbf{Z}[i]$ (anneau des entiers de $\mathbf{Q}[i]$).

La donnée d'un tel torseur $S$ permet de tordre
\[
\mathbb{E}\mathrm{l}_0 \overset{\mathrm{def}}{\simeq} \mathbf{C}/\mathbf{Z}[i]
\]
et trouver
\[
\mathbb{E}\mathrm{l}_0(S) \simeq \mathbb{E}\mathrm{l}_0 \wedge^{\mathbf{Z}/4\mathbf{Z}} S
\]
\note{à gauche, un croquis : le carré de sommets $1, i, -1, -i$ dans le
plan, les points du réseau $\mathbf{Z}[i]$ marqués, un carré plus petit
inscrit, centre $o$}

Si le carré est réalisé comme carré \add{euclidien}\note{« euclidien »
en interligne} dans un plan euclidien (canoniquement), de centre $o$,
\note{deux petits croquis : un carré biffé, puis un losange à
diagonales, étiqueté $V_S$} on en déduit dans le plan (vectoriel via
$o$) $V_S$ un syst.\ de deux \struck{axes de} coordonnées
\add{mod.\ signes} via les diagonales, d'où un \add{réseau}
\struck{\ill{}} $\Gamma_S$ dans $V_S$ (sous-groupe discret de rang 2) et
\[
\boxed{\mathbb{E}\mathrm{l}_0(S) =: X_S = V_S/\Gamma_S}\ .
\]
\note{la formule est encadrée ; à gauche, un croquis du losange dans un
quadrillage en pointillé}
Il y a dans $X_S$ un élément d'ordre 2 canonique $\eta_S$ (\uncertain{celui
fixé par automorphismes}), les deux autres él.\ d'ordre 2 sont échangés
par le générateur $i$ de
\[
\mu_4(\mathbf{C}) \simeq \mathrm{Aut}(X_S).
\]
\emph{\uncertain{CNB}} \ill{} \ill{} \ill{}\note{fin de ligne, au bas
de la page, illisible}

\page{6}
\note{p.~2 de l'auteur}
\ldots\ comme l'unique \uncertain{pt} \ill{} de ${}_2(X_S)$ qui
est fixe par $\mathrm{Aut}(X_S) \simeq \mathbf{Z}/4\mathbf{Z}$. En fait
\[
(X_S)^{i\text{-fix}} = \lbrace x \in X_S \mid ix = x \rbrace
= \lbrace e, \eta_S \rbrace .
\]
\note{l'exposant de $(X_S)$ est lu « $i$-fix » sans certitude ; il écrit
${}_2(X_S)$ pour le sous-groupe des points d'ordre 2}

La surface top.\ $X_S$ a une déc.\ cell.\ avec 1 sommet $e$, deux arêtes,
une face, \add{ces arêtes}\note{en interligne} correspondant aux deux
mailles du réseau $\Gamma_S$.

On considère $X_S$ comme muni de cette décomposition. Parfois c'est
l'image inverse de celle-ci par $2\,\mathrm{id}_{X_S}$ qui est plus
commode, 4 sommets (formant ${}_2(X_S)$), \struck{quatre} huit arêtes,
\struck{(NB \ill{})} 4 faces. On peut considérer
$\mathbf{Z}/4\mathbf{Z} \simeq \mu_4(\mathbf{C})$ comme le groupe
d'automorphismes commun de ces deux structures cellulaires.
\note{à gauche, un croquis : un carré hachuré découpé en quatre, sommets
$e$, $\eta'$, $\eta''$, $\eta$}

Le deuxième est mieux adapté pour la description du quotient
\[
\struck{\mathbb{E}\mathrm{l}_0(S)}\ X_S/\pm 1
\]
qui est une sphère, subdivisée en \struck{deux} \add{deux}
\uncertain{carrés}\note{au-dessus de « carrés », un mot lu « fus. » ;
le mot biffé devant est « deux »}, se raccordant suivant un équateur
\struck{carré} subdivisé en quatre arêtes par $e^{q}, \eta^{q},
\eta'^{q}, \eta''^{q}$.\note{l'exposant, lu $q$, est incertain. À gauche,
deux croquis de la sphère : l'un barré, l'autre avec l'équateur et les
quatre points $e^{q}, \eta^{q}, \eta'^{q}, \eta''^{q}$}

\page{8}
\note{p.~3 de l'auteur}
B) \emph{Courbes} \add{ell.}\ \emph{anharmoniques de type hexagonal.}

La catégorie de ces courbes équivalente à celle des hexagones orientés,
i.e.\ des torseurs $S$ sous
\[
\mathbf{Z}/6\mathbf{Z} \simeq \mu_6(\mathbf{C}) \simeq
\]
groupe des unités \add{alg.}\ de $\mathbf{Q}(\sqrt[3]{1})$
$\simeq$ groupe des unités \struck{de $\mathbf{Z}[\zeta] = \mathbf{Z}[T]$},
\[
\mathbf{Z}[\zeta] = \mathbf{Z}[T]/(1 + T + T^{2}) \qquad
(\zeta = \exp 2i\pi/3)
\]
\note{à gauche, un croquis : l'hexagone de sommets $\pm 1$, $\pm\zeta$,
$\pm\bar{\zeta}$, le réseau et ses points marqués}
\struck{On pose} On pose
\[
\mathbb{E}\mathrm{l}_1 \simeq \mathbf{C}/\mathbf{Z}[\zeta]
\]
\[
\mathbb{E}\mathrm{l}_1(S) \simeq \mathbb{E}\mathrm{l}_1
\wedge^{\mathbf{Z}/6\mathbf{Z}} S
\]
\note{dans la première formule, un argument biffé après $\mathbb{E}\mathrm{l}_1$}

\emph{Construction.} On réalise l'hexagone comb.\ $S$ euclidiennement
(canonique), dans \struck{\ill{}} \uncertain{vectoriel} euclidien $V_S$,
\add{origine en}\note{« origine » en interligne, au-dessus d'un mot
biffé} le barycentre de $S$. On désigne par $V_S$\note{sic ; la formule
qui suit appelle ce réseau $\Gamma_S$} le \add{réseau}\note{« réseau »
en interligne, au-dessus d'un mot biffé} engendré par $S$, et
\[
X_S = \mathbb{E}\mathrm{l}_1(S) \simeq V_S/\Gamma_S
\]
La structure euclidienne \emph{orientée} de $V_S$ définit la structure
complexe de $X_S$.

\page{10}
\note{p.~4 de l'auteur}
Le groupe $\mathbf{Z}/6\mathbf{Z}$ opère transitivement sur
$({}_2X_S) \setminus \lbrace 0 \rbrace$ (avec $\pm 1$ opérant
trivialement) --- donc ${}_2X_S \setminus \lbrace 0 \rbrace$ est un
torseur sous $(\mathbf{Z}/6\mathbf{Z})/\pm 1$
$\simeq \mathbf{Z}/3\mathbf{Z}$ (\struck{\ill{}} générateur correspondant
$= -\bar{\zeta}$, $\zeta = \exp\frac{2\pi}{6}$)\note{il écrit bien
$\exp\frac{2\pi}{6}$, sans $i$, et page 8 $\zeta = \exp 2i\pi/3$ ; le
$\bar{\zeta}$ est lu sans certitude}, \struck{\ill{} torseur} s'identifie
\[
\text{à} \quad S/\pm 1,
\]
en \uncertain{associant} à chaque $s \in S$ le milieu du segment
$[0,s]$, en \uncertain{prenant} son image dans $X_S$.

L'ens.\ des pts fixes de $\mathbf{Z}/3\mathbf{Z}$
\add{$\subset \mathbf{Z}/6\mathbf{Z}$}\note{en interligne} dans $X_S$ est
de card $3$, il s'obtient en prenant les \struck{\ill{}} barycentres des
six triangles de la subdivision centrale (bibarycentrique) de l'hexagone
(pts centrés) (tels que $(1-\zeta)/3$ etc.\ dans le cas épinglé), ils se
groupent en deux \uncertain{régiments} de $3$ mutuellement congrus mod
$\Gamma$, donc donnant deux éléments (symétriques) de $X_S$, d'ordre $3$,

\page{12}
\note{p.~5 de l'auteur}
qui avec l'origine $e$ forment l'ens.\ des pts fixes de $X_S$. (L'ens.\
des pts fixes sous $\mathbf{Z}/6\mathbf{Z}$ se réduit à $e$.) Les deux
pts fixes sous $\mathbf{Z}/3\mathbf{Z}$, distincts de $e$, correspondent
aux deux triangles inscrits \add{les deux}\note{en interligne} dans
l'hexagone $S$. (Triplets de côtés mutuellement non adjacents.)

\emph{Décompositions cellulaires canoniques.} Décomposition en deux faces
triangulaires, par trois arêtes, un seul sommet. Les faces triangulaires
correspondent aux éléments de $(X_S)^{\zeta} \setminus \lbrace e \rbrace$
(chaque face en contenant un et un seul) i.e.\ aux triangles inscrits
dans l'hexagone, les arêtes correspondent aux trois

\page{14}
\note{p.~6 de l'auteur}
diagonales \struck{de} de l'hexagone, ou encore aux trois pts d'ordre $2$
distincts de $e$.

\struck{\ill{}} Ici encore on peut « doubler » la décomposition
cellulaire, et on trouve une autre avec $4$ sommets, $12$ arêtes, et $8$
faces (triangulaires), qui permet de mieux comprendre le passage au
quotient par $\pm 1$, avec quatre sommets (dont un sommet privilégié,
correspondant à $e$), six arêtes \ill{} les valences $2 : 2$ : c'est la
structure tétraédrale \add{(orientée)}\note{en interligne}, avec un
sommet fixé, groupe d'automorphismes $\mathbf{Z}/3\mathbf{Z}$
(\struck{\ill{}} $\simeq \mathbf{Z}/6\mathbf{Z}/\pm 1$, OK !).
\note{à gauche, un croquis : l'hexagone subdivisé en triangles, un
losange de quatre petits triangles repassé en gras. Le « 12 » des arêtes
surcharge un autre chiffre. Le reste de la page est blanc ; la suite
paginée par lui s'arrête ici, à sa p.~6}

\section{$M_{0,4}$ \quad $M_{1,1}$}
\note{titre pris de la page 16, feuillet de papier listing vierge portant
seulement, de sa main, en haut à droite et l'un sous l'autre, « $M_{0\,4}$ »
et « $M_{1,1}$ » : couverture de ce qui suit, qui ne reçoit pas de numéro de page}

\page{17}
\note{papier listing, à l'encre noire, avec des croquis au crayon et aux
encres rouge, brune et verte ; pas de numéro de sa main}

4 sommets d'ordre 6 $+$ \struck{6} 6 sommets d'ordre 6 $+$ 12 sommets
d'ordre 3

24 triangles $+$ 4 hexagones

\note{sous ces deux lignes, deux croquis : un tétraèdre au crayon aux
sommets tronqués par de petits cercles ; à droite, un tétraèdre à l'encre
dont la subdivision barycentrique est tracée en rouge, avec, en marge :}
\marginal{mieux, cf.\ \uncertain{verso} (subdivision barycentrique)}

\[
s_0 = 4 + 6 + \struck{4}\,3 = 22 \qquad\struck{\ill{}}
\]
\[
s_1 = 12 + 4.\struck{9}\ \ldots = \struck{4}\,8\ \ldots
\qquad
s_2 = 4.7 = 28
\]
\note{les chiffres surchargés sont donnés tels qu'ils se lisent ; dans la
deuxième ligne le facteur et le résultat sont repassés et ne se
laissent pas lire avec certitude (« $4.9$ », « $48$ » ?)}

\note{à gauche, deux sphères au crayon et à l'encre : sur la première, un
fuseau hachuré en rouge, des points noirs et de petits hexagones de
points autour de trois sommets ; sur la seconde, un équateur et quatre
points. Au milieu, un petit croquis d'un « pantalon » à quatre trous.
Plus bas, un grand tétraèdre tronqué, arêtes vertes et brunes,
pointillés verts, dont un sommet est relié par une flèche à la note
de la marge droite :}
\marginal{ça fait $2^{4} - 1 = 15$ positions \uncertain{limites}
\uncertain{invariantes} \uncertain{possibles} : \ill{} \ill{} \ill{} de
vue !}\note{écrit en oblique ; lecture très incertaine}

\emph{4 types de points-base}
\begin{itemize}
  \item[$\bullet$] 3 structures carrées sur $I$ (ou $J$)
  \item[$\bullet$] 2 structures tétraédrales $=$ orientations
    \marginal{des faces}\note{point rouge}
  \item[$\bullet$] 6 structures \add{bi}pantalonniques (maillot ceinture)
    (carrée $+$ 1 diagonale) $\simeq$ repère de $J$
    \marginal{opposées}\note{point brun}
  \item[$\bullet$] 6 structures pantaloniques de genre $0$\note{lu « de
    genre gois » ; lecture incertaine} (quart de sorte) (\struck{\ill{}}
    carré, $+$ orientation de $I$ i.e.\ de $J$ $\simeq$ repère de $J$)
\end{itemize}
\note{à gauche de cette liste, en colonne :
$\omega_I \overset{\sim}{=} \omega_J$, $\mathrm{rep}(J)$, une double
flèche verticale, $J \times \omega_J$ ; une autre double flèche relie
le troisième et le quatrième point}

\emph{4 types de chemins élémentaires}$^{*}$
\begin{itemize}
  \item[a)] « quart de \uncertain{twist} » ($2\times3\times4$)
    (\uncertain{contour} d'un trou)
  \item[b)] « sixième de tour » ($2\times2\times3$) (sur équateur)
  \item[c)] pôle vers trou \uncertain{voisin} ($2\times2\times3$)
  \item[d)] pôle vers équateur \uncertain{voisin} ($2\times2\times3$)
\end{itemize}
2 types composés \uncertain{supplémentaires}.
\begin{itemize}
  \item[e)] trou vers trou ($2\times3$) ou « tiers de tour »
  \item[f)] demi \uncertain{twist} ($2\times2\times3$)
\end{itemize}
[et bien sûr, g) twist complet $2\times2\times3$]

\emph{relations} \ldots

$^{*}$ avec les deux sens de parcours, cela fait \emph{huit} types de
déformations : \uncertain{viendraient} a) deux pts d'un même
\uncertain{confluent}, b) deux pts d'un \uncertain{même} combinatoire
(ici \uncertain{encore}, deux pts d'un \uncertain{même} voisin)
\note{fin de page, écriture serrée ; la note est appelée par l'astérisque
du titre « 4 types de chemins élémentaires »}

\marginal{en outre les \uncertain{\ill{}}, on voit à la fois la
\uncertain{projection} du tétraèdre à la configuration carrée, et à la
configuration \uncertain{tétraédrale} \uncertain{pour la \ill{}} du genre
$0$ \ill{}}\note{bloc de texte à gauche, sous les sphères, en partie
illisible ; en bas à gauche, une sphère avec un équateur et quatre
points, deux méridiens en rouge}

\page{18}
\note{verso de la page 17, sur le dos d'un listing ; au crayon et à
l'encre. C'est le « verso » auquel renvoie la marge de la page 17}

\note{en haut à droite, une sphère avec un équateur et deux méridiens,
quatre points marqués sur l'équateur, des flèches le long des méridiens ;
à sa droite :}
\marginal{passage du carré aux deux tétraèdres}

Pour avoir le passage aux deux structures pantaloniques, il faut
introduire la \struck{\ill{}} subdivision barycentrique complète de la
carte carrée.

dans les deux cas \add{« spéciaux »}\note{en interligne, entre
guillemets} non digons (carré, tétraèdre) les positions adjacentes sont
toutes contenues dans la subdivision barycentrique

\note{en bas à gauche, un tétraèdre dont la subdivision barycentrique est
tracée, un chemin repassé en gras ; à sa droite :}
cette carte (subdivision bar.\ du \add{\uncertain{nouveau}}\note{en
interligne, au-dessus de « tétraèdre »} tétraèdre) contient également les
subdivisions barycentriques des trois cartes « carrés » associées, (dont
l'une est marquée en pointillés rouges)

\note{au bas de la page, une sphère au crayon avec équateur, méridiens et
quatre points marqués}

\page{19}
\note{dos d'un listing, au crayon}

\note{en haut, une suite de croquis de sphères : à gauche une sphère à
équateur avec quatre flèches ; au centre une sphère dont l'équateur porte
les points $a, b$ et, au-dessus, $c, d$, d'où partent quatre flèches :
vers le haut, marquée $\rho$, deux sphères et un tétraèdre ; vers la
gauche, marquée $\sigma$, une sphère méridienne et deux sphères accolées,
avec la lettre $\ell$ au-dessous ; vers la droite, deux sphères ; vers le
bas, deux sphères et un tétraèdre}

\[
I = S_{Q} = \coprod_{d \in D_Q} d \qquad
A_{Q} = \prod_{d \in D_Q} d \qquad
C_{Q} = \bigwedge_{d \in D_Q} d
\]
\[
\omega_{Q} = C_{Q} \wedge D_{Q} \qquad
\omega_{I} \simeq \omega_{J} \simeq C_{Q} \simeq J \setminus \lbrace Q \rbrace
\]
\note{la première égalité, $I = S_Q$, surcharge une lettre ; les signes
$\simeq$ de la dernière ligne sont repassés}

\[
\begin{array}{l}
\quad a\ b\ c\ d \\
(a\,d)(b\,c) = Q_0 \\
(b\,d)(a\,c) = Q_1 \\
(c\,d)(a\,b) = Q_{\infty}
\end{array}
\]
\note{ces trois lignes sont encadrées à gauche et en haut}

\[
(\text{tétraèdre \emph{orienté} } + \text{diagonale})
\longleftrightarrow (\text{carré} + \text{codiagonale})
\]
\note{sous la double flèche, une accolade : « $\omega_I \times J \simeq J
\times \omega_I$ » ; sous le tétraèdre, un $\omega_I \times J$ biffé}

\[
0 \to V(J) \to \mathfrak{S}_{I} \to \mathfrak{S}_{J} \to 1
\qquad\qquad
\boxed{\begin{cases} \rho^{3} = \sigma^{2} = 1 \\ \rho\sigma = \ell \end{cases}}
\]

\begin{tikzcd}[column sep=small, row sep=small]
1 \arrow[r] & \mathcal{T}^{!}_{0,4} \arrow[r] \arrow[d, no head, "\wr"'] & \mathcal{T}_{0\,4} \arrow[r] \arrow[d] & \mathfrak{S}_4 \arrow[r] \arrow[d] & 1 \\
1 \arrow[r] & \pi_{0,3} \arrow[r] & \mathcal{T}'_{0\,4} \arrow[r] & \mathfrak{S}_3 \arrow[r] & 1 \\
1 \arrow[r] & \pi_{0,3}(2,\mathbf{Z}) \arrow[r] \arrow[u, no head, "\wr"] & \mathcal{T}_{1,1} \arrow[r] \arrow[u] & \widetilde{\mathfrak{S}}_3 \arrow[u]
\end{tikzcd}

\note{le diagramme est au crayon. Annotations : $\mathfrak{S}_4 \simeq
\mathrm{Aff}(2, \mathbf{Z}/2\mathbf{Z})$ ; $\mathfrak{S}_3 \simeq
\mathrm{Sl}(2, \mathbf{Z}/2\mathbf{Z})$ ; $\widetilde{\mathfrak{S}}_3 \simeq
\mathrm{Sl}(2, \mathbf{Z}/4\mathbf{Z})/\pm$ (le signe final est lu
« $/-$ ») ; $\mathcal{T}_{1,1} \simeq \mathrm{Sl}(2,\mathbf{Z})$ ; la
flèche $\mathcal{T}_{0\,4} \to \mathfrak{S}_4$ part d'un exposant
griffonné illisible sur $\mathcal{T}_{0\,4}$, et « \uncertain{cart} » est
écrit près du carré de droite. Le premier terme de la dernière ligne est
lu $\pi_{0,3}(2,\mathbf{Z})$ sans certitude : il semble écrire
$\pi_{0,3}$ surchargé d'un « $(2,\mathbf{Z})$ » plus pâle ; la première
ligne commence par un signe \ill{} avant la flèche}

\page{20}
\note{dos d'un listing ; au crayon, sans texte : un tétraèdre, entouré de
six traits rayonnants, et à droite un petit octaèdre (ou deux tétraèdres
croisés) traversé par un axe}

\end{document}
